\scanpage{176}

Lorsque \deleted{$X$ est} \deleted{\unsure{l'on est dans les conditions de}} les anneaux localisés stricts de $X$ sont factoriels, on tire aisément de (5.3.)
\[
H^1(X,\uline{\mathrm{Div}}_X)=0
\]
d'où \deleted{la proposition}

\uline{Proposition} 5.4. Supposons que l'on soit dans les conditions de 5.1. a), \inserted{$X$ étant irréductible de pt gén. $y$,} Alors les homomorphismes canoniques
\diagram{\begin{tikzcd}[column sep=large]
H^2(X,G_{mX})\arrow[r,no head]&H^2(X,R_X^*)\arrow[r]&H^2(y,G_{my})=\operatorname{Br}(y)
\end{tikzcd}}
sont injectifs.

Comme $\operatorname{Br}(X)\subset H^2(X,G_{mX})$, il s'ensuit

\uline{Corollaire} 5.5. Sous les $=^{s}$ conditions,
\diagram{\begin{tikzcd}[column sep=large]
\operatorname{Br}(X)\arrow[r,no head]&\operatorname{Br}(y)
\end{tikzcd}}
est injectif.

Ainsi, $\operatorname{Br}(X)$ s'identifie \inserted{alors} à un sous-groupe du groupe de Brauer $\operatorname{Br}(K)$ du corps des fonctions $K=k(y)$ de $X$.

\uline{Théorème} 5.6. \inserted{Soit $X$ un préschéma, $\xi\in H^2(X,G_m)$}
a) Si $X$ est un schéma local, et \deleted{$\xi\in H^2(X,G_{mX})$} pour que $\xi$ \deleted{provienne d'un élément de} \inserted{soit dans $\operatorname{Br}(X)$, il} f et s. qu'il existe un\deleted{e extension} morphisme fini étale surjectif (galoisien si on veut) $f:X'\longrightarrow X$, \deleted{tel que $\xi|X'\to0$} qui \deleted{splitte $\xi$} annule $\xi$.

\margin{175}
